% ================= SEKCJA 9 =================
\section{Wykorzystanie narzędzi AI w procesie realizacji projektu}

W trakcie realizacji projektu \texttt{WorkFlow} autor korzystał z narzędzi opartych na sztucznej inteligencji jako wsparcia w analizie problemów, porządkowaniu zadań, redagowaniu dokumentacji oraz weryfikacji wybranych rozwiązań technicznych. Wykorzystywane były przede wszystkim narzędzia ChatGPT oraz Gemini, przy czym ChatGPT był używany częściej ze względu na lepsze dopasowanie odpowiedzi do kontekstu projektu i możliwość prowadzenia dłuższej, spójnej pracy nad kolejnymi etapami systemu.

Narzędzia AI nie zastępowały samodzielnej realizacji projektu. Pełniły funkcję konsultacyjną, podobną do korzystania z dokumentacji technicznej, wyszukiwarki, forów programistycznych lub narzędzi wspierających analizę kodu. Ostateczne decyzje projektowe, implementacyjne i redakcyjne były podejmowane przez autora projektu, a zaproponowane rozwiązania były każdorazowo weryfikowane w praktyce.

\subsection{Zakres wykorzystania narzędzi AI}

Narzędzia AI były wykorzystywane głównie w tych obszarach, w których przyspieszały analizę lub pomagały uporządkować pracę. Dotyczyło to zarówno części programistycznej, jak i dokumentacyjnej projektu.

Najważniejsze zastosowania obejmowały:
\begin{itemize}
    \item analizę błędów pojawiających się podczas pracy nad frontendem i backendem,
    \item interpretację komunikatów TypeScript, ESLint, Prisma, Docker oraz Node.js,
    \item porządkowanie kolejności zadań wykonywanych w projekcie,
    \item przygotowywanie propozycji komend testowych i diagnostycznych,
    \item analizę problemów z uruchomieniem aplikacji w kontenerach Docker,
    \item wsparcie przy konfiguracji środowiska produkcyjnego,
    \item pomoc w przygotowaniu danych demonstracyjnych,
    \item redagowanie i aktualizację dokumentacji technicznej,
    \item upraszczanie zbyt rozbudowanych lub nieprecyzyjnych opisów.
\end{itemize}

Szczególnie istotne było wykorzystanie AI jako narzędzia do bieżącej analizy problemów. W przypadku pojawienia się błędu autor mógł przedstawić komunikat z terminala lub fragment kodu, a następnie otrzymać propozycję przyczyny problemu oraz możliwych kroków diagnostycznych.

\subsection{Wsparcie podczas implementacji}

Podczas implementacji systemu narzędzia AI były wykorzystywane do analizy i wyjaśniania problemów technicznych. Dotyczyło to między innymi błędów kompilacji TypeScript, ostrzeżeń ESLint, problemów z działaniem formularzy, konfiguracją zmiennych środowiskowych oraz uruchomieniem aplikacji w Dockerze.

Przykładowe obszary, w których AI wspierało pracę implementacyjną, to:
\begin{itemize}
    \item analiza błędów związanych z typami danych w TypeScript,
    \item poprawianie problemów wykrywanych przez ESLint,
    \item diagnozowanie problemów z logowaniem użytkownika,
    \item analiza konfiguracji frontendowego adresu API,
    \item pomoc przy konfiguracji Dockerfile dla frontendu i backendu,
    \item analiza błędów uruchamiania aplikacji w kontenerach,
    \item wsparcie przy pracy z Prisma i inicjalizacji bazy danych,
    \item przygotowanie oraz weryfikacja seedów bazy danych.
\end{itemize}

W praktyce AI nie było traktowane jako źródło gotowych, bezwarunkowo poprawnych rozwiązań. Każda propozycja wymagała sprawdzenia w kodzie, uruchomienia aplikacji oraz oceny, czy pasuje do aktualnej struktury projektu. W wielu przypadkach wygenerowane sugestie wymagały modyfikacji, uproszczenia albo dostosowania do istniejących plików i konwencji.

\subsection{Wsparcie przy konfiguracji środowiska}

Istotnym elementem projektu było przygotowanie środowiska uruchomieniowego opartego na Docker Compose. W tym obszarze narzędzia AI pomagały w analizie zależności pomiędzy kontenerami, sposobie przekazywania zmiennych środowiskowych oraz diagnozowaniu błędów wynikających z różnic pomiędzy środowiskiem lokalnym a kontenerowym.

AI wspierało między innymi analizę problemów związanych z:
\begin{itemize}
    \item brakiem zbudowanych plików aplikacji w katalogu \texttt{dist},
    \item błędną ścieżką startową backendu,
    \item brakiem wymaganych zmiennych środowiskowych w kontenerze,
    \item konfiguracją połączenia Prisma z bazą PostgreSQL,
    \item przekierowaniem ruchu przez Nginx,
    \item różnicą pomiędzy adresem API widocznym dla przeglądarki a adresem API widocznym wewnątrz sieci Docker.
\end{itemize}

W tym zakresie AI pełniło rolę pomocniczą przy interpretacji logów i wskazywaniu kolejnych komend diagnostycznych. Ostateczna weryfikacja odbywała się poprzez uruchomienie kontenerów, analizę logów, testowanie endpointów za pomocą \texttt{curl} oraz sprawdzenie działania aplikacji w przeglądarce.

\subsection{Wsparcie przy danych demonstracyjnych}

Narzędzia AI zostały również wykorzystane podczas przygotowywania danych demonstracyjnych. Dane te były potrzebne do prezentacji działania systemu bez konieczności ręcznego wprowadzania dużej liczby rekordów.

Wsparcie AI obejmowało:
\begin{itemize}
    \item zaplanowanie zakresu danych demonstracyjnych,
    \item przygotowanie struktury seeda prezentacyjnego,
    \item analizę błędów TypeScript w skrypcie seedującym,
    \item uporządkowanie danych testowych tak, aby obejmowały najważniejsze moduły systemu,
    \item dodanie zabezpieczenia przed przypadkowym uruchomieniem destrukcyjnego seeda.
\end{itemize}

Dzięki temu możliwe było przygotowanie danych obejmujących użytkowników, zakłady, pracowników, projekty, czytniki RCP, umowy, certyfikaty, dokumenty, nieobecności, wpisy czasu pracy oraz przykładowy eksport rozliczeniowy. Dane te zostały następnie sprawdzone w działającej aplikacji.

\subsection{Wsparcie przy dokumentacji}

Drugim ważnym obszarem wykorzystania narzędzi AI była dokumentacja projektu. W tym przypadku AI pomagało przede wszystkim w porządkowaniu treści, aktualizacji opisów do aktualnego stanu aplikacji oraz usuwaniu zbyt ogólnych lub przesadnie marketingowych sformułowań.

W dokumentacji AI wspierało:
\begin{itemize}
    \item redakcję opisu celu projektu,
    \item uporządkowanie wymagań funkcjonalnych i niefunkcyjnych,
    \item aktualizację scenariuszy użycia,
    \item opis modelu danych i diagramu ERD,
    \item opis stosu technologicznego,
    \item przygotowanie instrukcji prezentacji projektu,
    \item ujednolicenie stylu językowego dokumentacji.
\end{itemize}

Ważnym założeniem było zachowanie formalnego charakteru dokumentacji bez nadmiernego rozbudowywania opisów. Z tego powodu część propozycji była upraszczana lub odrzucana, jeżeli brzmiała zbyt marketingowo albo sugerowała funkcje, które nie zostały zaimplementowane w aktualnej wersji projektu.

\subsection{Weryfikacja odpowiedzi generowanych przez AI}

Korzystanie z narzędzi AI wymagało zachowania ostrożności. Odpowiedzi generowane przez modele językowe mogą być niepełne, nieaktualne albo niedopasowane do konkretnego projektu. Z tego względu żadna propozycja nie była traktowana jako automatycznie poprawna.

Weryfikacja odbywała się poprzez:
\begin{itemize}
    \item analizę istniejącego kodu źródłowego,
    \item uruchamianie aplikacji lokalnie,
    \item sprawdzanie logów backendu, frontendu i kontenerów,
    \item testowanie endpointów za pomocą zapytań HTTP,
    \item kontrolę działania funkcji w interfejsie użytkownika,
    \item sprawdzanie zmian w systemie kontroli wersji Git,
    \item porównywanie dokumentacji z rzeczywistym zakresem aplikacji.
\end{itemize}

Taki sposób pracy ograniczał ryzyko bezrefleksyjnego przyjęcia błędnych sugestii. AI było wykorzystywane jako narzędzie wspomagające analizę i organizację pracy, a nie jako samodzielny wykonawca projektu.

\subsection{Porównanie wykorzystanych narzędzi}

W trakcie pracy wykorzystywane były dwa narzędzia: Gemini oraz ChatGPT. Oba rozwiązania mogły wspierać analizę problemów i redakcję treści, jednak w praktyce większe znaczenie w projekcie miało narzędzie ChatGPT.

Gemini było wykorzystywane pomocniczo, głównie do porównywania propozycji odpowiedzi lub uzyskania alternatywnego spojrzenia na dany problem. ChatGPT okazał się bardziej przydatny podczas dłuższej pracy nad projektem, ponieważ lepiej utrzymywał kontekst kolejnych etapów, takich jak rozwój modułów, poprawki błędów, konfiguracja Docker oraz aktualizacja dokumentacji.

Nie oznacza to jednak, że odpowiedzi ChatGPT były przyjmowane bez weryfikacji. Narzędzie to było traktowane jako asystent konsultacyjny, którego propozycje należało sprawdzić i dostosować do rzeczywistego kodu projektu.

\subsection{Wpływ AI na przebieg pracy}

Wykorzystanie narzędzi AI wpłynęło przede wszystkim na tempo pracy i sposób organizacji zadań. Dzięki możliwości szybkiej analizy błędów i propozycji kolejnych kroków łatwiej było przechodzić przez problemy techniczne oraz utrzymywać porządek w realizowanych zadaniach.

Największe korzyści z wykorzystania AI to:
\begin{itemize}
    \item szybsza analiza błędów i komunikatów diagnostycznych,
    \item łatwiejsze planowanie kolejnych etapów pracy,
    \item możliwość szybkiego porównania kilku wariantów rozwiązania,
    \item wsparcie przy redakcji dokumentacji,
    \item większa konsekwencja w nazewnictwie i opisie modułów,
    \item pomoc w przygotowaniu komend testowych i diagnostycznych.
\end{itemize}

Jednocześnie wykorzystanie AI nie zwalniało autora z odpowiedzialności za projekt. Kod musiał zostać uruchomiony, przetestowany i dopasowany do istniejącej struktury aplikacji. Dokumentacja musiała zostać porównana z rzeczywistym zakresem systemu, aby nie opisywała funkcji, które nie są częścią projektu.

\subsection{Podsumowanie}

Narzędzia AI były istotnym wsparciem podczas realizacji projektu \texttt{WorkFlow}, szczególnie w zakresie analizy błędów, organizacji pracy i redakcji dokumentacji. Ich rola miała jednak charakter pomocniczy. Projekt, jego zakres, decyzje techniczne, testowanie oraz ostateczny kształt dokumentacji pozostawały pod kontrolą autora.

Takie wykorzystanie AI można traktować jako element współczesnego procesu pracy programistycznej. Narzędzia te pomagają szybciej analizować problemy i porządkować informacje, ale wymagają krytycznego podejścia, weryfikacji oraz świadomego dostosowania wygenerowanych sugestii do rzeczywistego projektu.